%%%%%%%%%%  scan idx 358  |  folio 345  %%%%%%%%%%
\[
  (1+\widehat{\cdot}{}^{+})\times(1+\widehat{\cdot}\cdot{}^{+})\xrightarrow{\ \ast\ }1+\widehat{\cdot}{}^{+},
\]
\[
  1+\widehat{\cdot}{}^{+}\xrightarrow{\ a\ }1+\widehat{\cdot}{}^{+},
\]
définissant sur le foncteur \(1+\widehat{\cdot}{}^{+}\) une multiplication et des opérations de \(K\). D'après 1.13, la multiplication est définie par sa valeur sur l'élément \((1+X,1+Y)\) de \((1+\widehat{K[X,Y]}{}^{+})^{2}\)~; elle sera notée \(\ast\)~; on a donc :
\begin{equation}\tag{6.1.1}
  (1+X)\ast(1+Y)=\frac{1+X+Y}{(1+X)(1+Y)}.
\end{equation}
Le \(n\)-ième polynôme universel de cette loi est la composante isobare de poids \(n\) de \(P'_n\) (1.15) et sera noté \(P''_n\).

Les opérations de \(K\) seront encore notées par l'exponentielle (3.8). Les polynômes universels définissant ces opérations sont les composantes isobares de plus bas degré des \(R_{n,a}\) (1.15)~; ceux-ci étant isobares, ce sont les \(R_{n,a}\) eux-mêmes.

Soit \(A\in\operatorname{Ob}(\mathcal G)\)~; l'anneau \(1+\widehat A^{+}\) défini par les lois précédentes est en général sans élément unité~; on étend donc cet anneau par adjonction des scalaires de \(K\), en introduisant le nouvel anneau \(K\times(1+\widehat A^{+})\), muni des lois de composition définie par : \(\forall\,a,a'\in A\), \(\varphi,\varphi'\in1+\widehat A^{+}\),
\[
  (a,\varphi)+(a',\varphi')=(a+a',\varphi\varphi')~;
\]
\begin{equation}\tag{6.1.2}
  (a,\varphi).(a',\varphi')=(aa',\varphi^{a'}\cdot{\varphi'}^{\,a}\cdot(\varphi\ast\varphi')).
\end{equation}

L'anneau ainsi construit admet (1.1) pour l'élément unité, et \(1+\widehat A^{+}\) s'identifie à l'idéal \((0)\times(1+\widehat A^{+})\). Il est muni de façon naturelle d'une structure de \(K\)-algèbre, telle que :
\begin{equation}\tag{6.1.3}
  a.(a',\varphi)=(a,1).(a',\varphi)=(aa',\varphi^{a}).
\end{equation}

